\subsection{Type Soundness}
\label{SoundnessSection}
The value of a type soundness theorem for Remora is not only assurance
that well-typed programs do not suffer from shape-mismatching errors.
It also ensures that the types ascribed to program terms
accurately describe the shapes of the data those terms compute.
That is the guarantee that justifies a compiler's use of the type system
as a static analysis for array shape.

With supporting lemmas,
such as canonical forms and substitution,
already taken care of,
we now establish progress and preservation lemmas.
Since we have not committed to
a collection of primitive operators that are all total functions,
the progress lemma acknowledges the possibility of non-shape errors,
such as division by zero.
However, we do assume that any value returned by a primitive operator
inhabits that operator's output type.

\newcommand{\misapp}{\app{\arrlit{\primop}{}}{\sequence{\val}}}
\begin{lemma}[Progress]
  \label{Progress}
  Given an expression $\expr$ such that
  $\typeof{\emptyEnv}{\emptyEnv}{\emptyEnv}{\expr}{\type}$,
  one of the following holds:
  \begin{itemize}
  \item{$\expr$ is a value $\val$}
  \item{There exists $\expr'$ such that $\expr \mapsto \expr'$}
  \item{$\expr$ is
    $\ctxt\left[\misapp\right]$
    where $\primop$ is a partial function applied to
    appropriately-typed values outside its domain.}
  \end{itemize}
\end{lemma}
\begin{sproof}[We use induction on the derivation of
  $\typeof{\emptyEnv}{\emptyEnv}{\emptyEnv}{\expr}{\type}$.
  We consider only cases for typing rules which apply to expressions (as opposed to atoms).
  Since we do not reduce under a binder,
  our assumed type derivation ensures
  that the reducible subexpression of $\expr$
  is also typable using an empty environment.

  An {\tt array} form which is not already a value
  must have some non-value atom.
  That atom must itself contain a non-value expression,
  with its own type derivation.
  So the induction hypothesis implies that
  it can take a reduction step or is a mis-applied primitive operator.
  Similar reasoning applies to {\tt frame} forms:
  either we have a {\it collapse} redex,
  or some cell subexpression in the frame can make progress.

  An {\tt unbox} form can either
  make progress in the box position (via the induction hypothesis)
  or take an {\it unbox} step.
  Similarly, a type or index applications can make progress in function position
  or take $\mathit{t\beta}$ or $\mathit{i\beta}$ step.

  The function application case splits into subcases depending on
  whether the function and argument arrays are fully reduced
  and if so what their frame shapes are.
  If they are all value forms,
  we have all scalar frames (a $\beta$ or $\delta$ redex)
  or all identical non-scalar frames (a {\it map} redex),
  or non-identical prefix-compatible frames (a {\it lift} redex).
  Prefix-incompatible frames are ruled out by the type derivation.]
  We use induction on the derivation of
  $\typeof{\emptyEnv}{\emptyEnv}{\emptyEnv}{\expr}{\type}$.
  We consider only cases for typing rules which apply to expressions (as opposed to atoms).
  Note that a {\sc T-Var} derivation is impossible with an empty type environment.
  \paragraph{Case {\sc Array}}
  \[{\infr[T-Array]
    {\sequence{\typeof{\emptyEnv}{\emptyEnv}{\emptyEnv}{\atom}{\type_a}}
      \qquad
      \kindof{\emptyEnv}{\emptyEnv}{\type_a}{\kindatom}
      \qquad
      \mathit{Length}\llb \sequence{\atom} \rrb = \prod{\sequence{\nat}}}
    {\typeof{\emptyEnv}{\emptyEnv}{\emptyEnv}
      {\arrlit{\sequence{\atom}}{\sequence{\nat}}}
      {\typearray{\type_a}{\idxshape{\sequence{\nat}}}}}}\]
  We have two possibilities.
  One is that all of $\sequence{\atom}$ are atomic values, $\sequence{\atval}$,
  in which case $e = \typearray{\sequence{\atval}}{\sequence{\nat}}$ is a value.
  Otherwise, at least one $\atom$ is not an atomic value.
  The non-value must be of the form
  $\atom_s = \dsum{\sequence{\notevar{\var}{\sort}}}{\expr_s}{\type_a}$,
  with non-value $\expr_s$.
  The derivation of $\typeof{\emptyEnv}{\emptyEnv}{\emptyEnv}{\atom}{\type_a}$
  implies via the induction hypothesis that
  either there exists some $\expr_s'$ such that $\expr_s \mapsto \expr_s'$
  or $\expr_s$ has the form $\ctxt_s\left[\misapp\right]$
  with misapplied partial function $\primop$.
  We build the evaluation context
  $\ctxt = \arrlit{\sequence{\atom_0} \;
    \dsum{\sequence{\notevar{\var}{\sort}}}{\ctxt_s}{\type_a}
    \; \sequence{\atom_1}}
  {\sequence{\nat}}$
  where $\sequence{\atom_0}$ and $\sequence{\atom_1}$
  are the atoms appearing respectively before and after $\atom_s$ in $\expr$.
  If $\expr_s \mapsto \expr_s'$, then $\ctxt\left[\expr_s\right] \mapsto \ctxt\left[\expr_s'\right]$.
  Otherwise, $\expr = \ctxt\left[\misapp\right]$.

  \paragraph{Case {\sc Frame}}
  \[{\infr[T-Frame]
    {\sequence{\typeof{\emptyEnv}{\emptyEnv}{\emptyEnv}
        {\expr_a}
        {\typearray{\type_a}{\idx_a}}}
      \qquad
    \kindof{\emptyEnv}{\emptyEnv}{\typearray{\type_a}{\idx_a}}{\kindarray}
    \\\\
    \mathit{Length}\llb \sequence{\expr} \rrb = \prod{\sequence{\nat}}}
    {\typeof{\emptyEnv}{\emptyEnv}{\emptyEnv}
      {\frm{\sequence{\expr_a}}{\sequence{\nat}}}
      {\typearray{\type_a}{\idxappend{\idxshape{\sequence{\nat}} \; \idx_a}}}}}\]
  By the induction hypothesis,
  each of $\sequence{\expr_a}$ is a value, is reducible,
  or is a primitive operator misapplication
  with the form $\ctxt_a\left[\misapp\right]$.
  If they are all values, then
  each has the form $\arrlit{\sequence{\atval}}{\sequence{\nat'}}$, so
  $\expr \mapsto$
  {\tt (array ($\sequence{\nat}$ $\sequence{\nat'}$)
    $\mathit{Concat}\llbracket(\atval \sequence{})$ $\sequence{}$
    $\rrbracket$)}.
  % $\arrlit
  % {\mathit{Concat}\llb\sequence{\parens{\sequence{\atval}}}\rrb}
  % {\sequence{\nat}\;\sequence{\nat'}}$.
  Note that all of the array literals serving as cells in the frame must have the same shape,
  or their types would differ, making $\expr$ ill-typed.

  If some $\expr_e \in \sequence{\expr_a}$ is not a value,
  then the induction hypothesis implies that
  it is reducible to $\expr_e'$
  or it has the form $\ctxt_e\left[\misapp\right]$.
  We construct the evaluation context
  $\ctxt =$
  {\tt (frame ($\nat$ $\sequence{}$)
    $\expr_0$ $\sequence{}$ $\expr_e$ $\expr_1$ $\sequence{}$)},
  % $\ctxt = \frm{\sequence{\expr_0} \; \expr_e \; \sequence{\expr_1}}{\sequence{\nat}}$,
  where $\parens{\sequence{\expr_a}} = \parens{\sequence{\expr_0} \; \expr_e \; \sequence{\expr_1}}$.
  If $\expr_e \mapsto \expr_e'$, then $\ctxt\left[\expr_e\right] \mapsto \ctxt\left[\expr_e'\right]$.
  Otherwise, $\expr = \ctxt\left[\misapp\right]$.

  \paragraph{Case {\sc TApp}}
  \[{\infr[T-TApp]
    {\typeof{\emptyEnv}{\emptyEnv}{\emptyEnv}
      {\expr_f}
      {\typearray{\typeuniv{\sequence{\notevar{\var}{\kind}}}{\typearray{\type_u}{\idx_u}}}{\idx_f}}
      \\\\
    \sequence{\kindof{\emptyEnv}{\emptyEnv}{\type_a}{\kind}}}
    {\typeof{\emptyEnv}{\emptyEnv}{\emptyEnv}
      {\tapp{\expr_f}{\sequence{\type_a}}}
      {\typearray{\seqsubst{\type_u}{\var}{\type_a}}{\idxappend{\idx_f \; \idx_u}}}}}\]
  By the induction hypothesis,
  $\expr_f$ is a value, is reducible, or misapplies a partial function.
  If $\expr_f$ is a value of type
  ${\typearray{\typeuniv{\sequence{\notevar{\var}{\kind}}}{\typearray{\type_u}{\idx_u}}}{\idx_f}}$,
  then Lemma \ref{CFArray} (canonical forms) implies
  that it is an array literal containing type abstractions---%
  \ie, $\expr_f = \arrlit{\sequence{{\tlam{\sequence{\notevar{\var_u}{\kind}}}{\val_u}}}}{\sequence{\nat}}$.
  This means we have a $\mathit{t\beta}$ redex.%,
  % which reduces to $\frm{\sequence{\seqsubst{\val_u}{\var_u}{\type_a}}}{\sequence{\nat}}$.
  If $\expr_f$ is itself reducible to $\expr_f'$,
  then for some context $\ctxt_f$ and redex $\expr_r$,
  $\expr_f = \ctxt_f\left[\expr_r\right]$, and $\expr_r \mapsto \expr_r'$.
  % So $\ctxt = \tapp{\ctxt_f}{\sequence{\type_a}}$,
  % thus $\expr = \ctxt\left[\expr_r\right] \mapsto \ctxt\left[\expr_r'\right]$.
  Alternatively, $\expr_f = \ctxt_f\left[\misapp\right]$.
  Either way, construct $\ctxt = \tapp{\ctxt_f}{\sequence{\type_a}}$.
  In the first case, $\expr = \ctxt\left[\expr_r\right] \mapsto \ctxt\left[\expr_r'\right]$.
  In the second, $\expr = \ctxt\left[\misapp\right]$.
  \paragraph{Case {\sc IApp}}
  \[{\infr[T-IApp]
    {\typeof{\emptyEnv}{\emptyEnv}{\emptyEnv}
      {\expr_f}
      {\typearray{\typedprod{\sequence{\notevar{\var}{\sort}}}{\typearray{\type_p}{\idx_p}}}{\idx_f}}
      \qquad
      \sequence{\sortof{\emptyEnv}{\idx_a}{\sort}}}
    {\typeof{\emptyEnv}{\emptyEnv}{\emptyEnv}
      {\iapp{\expr_f}{\sequence{\idx_a}}}
      {\typearray{\seqsubst{\type_p}{\var}{\idx_a}}{\idxappend{\idx_f \; \seqsubst{\idx_p}{\var}{\idx_a}}}}}}\]
  As in the previous case,
  $\expr_f$ is a value, is reducible, or misapplies a partial function.
  If $\expr_f$ is a value, the canonical forms lemma implies that it has the form
  {\tt
    (array ($\nat$ $\sequence{}$) (I\cl (($\var_p$ $\sort$) $\sequence{}$) $\val_p$) $\sequence{}$)},
  which makes $\expr$ an $\mathit{i\beta}$ redex.
  Otherwise $\expr_f$ itself is reducible,
  \ie, of the form $\ctxt_f\left[\expr_r\right]$ for some rede $\expr_r$ reducing to $\expr_r'$,
  \emph{or} it is of the form $\ctxt_f\left[\misapp\right]$.
  Let the context $\ctxt = \iapp{\ctxt_f}{\sequence{\idx_a}}$.
  Then we have either $\expr = \ctxt\left[\expr_r\right] \mapsto \ctxt\left[\expr_r'\right]$
  or $\expr = \ctxt\left[\misapp\right]$.

  \paragraph{Case {\sc Unbox}}
  \[{\infr[T-Unbox]
    {\typeof{\emptyEnv}{\emptyEnv}{\emptyEnv}
      {\expr_s}
      {\typearray{\typedsum{\sequence{\notevar{\var_i'}{\sort}}}{\type_s}}{\idx_s}}
      \\\\
      \typeof
      {\sequence{\hassort{\var_i}{\sort}}}{\emptyEnv}{\hastype{\var_e}{\seqsubst{\type_s}{\var_i'}{\var_i}}}
      {\expr_b}
      {\typearray{\type_b}{\idx_b}}
      \\\\
      \kindof{\emptyEnv}{\emptyEnv}{\typearray{\type_b}{\idx_b}}{\kindarray}}
    {\typeof{\emptyEnv}{\emptyEnv}{\emptyEnv}
      {\dproj{\sequence{\var_i}}{\var_e}{\expr_s}{\expr_b}}
      {\typearray{\type_b}{\idxappend{\idx_s \; \idx_b}}}}}\]
  By the induction hypothesis,
  $\expr_s$ is a value, reducible, or a misapplication of a partial function.
  If $\expr_s$ is a value, the canonical forms lemma implies
  $\expr_s = \arrlit{\dsum{\sequence{\idx_s}}{\val}{\type}}{}$,
  so $\expr$ is an $\mathit{unbox}$ redex.
  If $\expr_s$ is reducible,
  \ie, it is $\ctxt_s\left[\expr_r\right]$ where the redex $\expr_r \mapsto \expr_r'$,
  then let $\ctxt = \dproj{\sequence{\var_i}}{\var_e}{\ctxt_s}{\expr_b}$.
  Thus $\expr = \ctxt\left[\expr_r\right] \mapsto \ctxt\left[\expr_r'\right]$.
  Otherwise, $\expr_s =$
  $\ctxt_s[$
  {\tt ((array () $\primop$) $\val$ $\sequence{}$)}
  $]$,
  so the same construction of $\ctxt$ gives
  $\expr = \ctxt\left[\misapp\right]$.
  \paragraph{Case {\sc App}}
  \[{\infr[T-App]
    {\typeof{\emptyEnv}{\emptyEnv}{\emptyEnv}{\expr_f}
      {\typearray
        {\typefun
          {\sequence{\typearray{\type_i}{\idx_i}}}
          {\typearray{\type_o}{\idx_o}}}
        {\idx_f}}
      \\\\
    \sequence{\typeof{\emptyEnv}{\emptyEnv}{\emptyEnv}
      {\expr_a}
      {\typearray{\type_i}{\idxappend{\idx_a \; \idx_i}}}}
    \qquad
    \idx_p = \mathit{Max}\llb \idx_f \; \sequence{\idx_a} \rrb}
    {\typeof{\emptyEnv}{\emptyEnv}{\emptyEnv}
      {\app{\expr_f}{\sequence{\expr_a}}}
      {\typearray{\type_o}{\idxappend{\idx_p \; \idx_o}}}}}\]
  Suppose $\expr_f$ is not a value.
  Then the induction hypothesis implies that
  \emph{either} $\expr_f \mapsto \expr_f'$,
  meaning $\expr_f = \ctxt_f\left[\expr_r\right]$
  for some redex $\expr_r \mapsto \expr_r'$
  \emph{or} $\expr_f = \ctxt_f\left[\misapp\right]$
  with incompatible but properly-typed arguments.
  Let $\ctxt = \app{\ctxt_f}{\sequence{\expr_a}}$.
  If $\expr_f \mapsto \expr_f'$, with $\expr_f = \ctxt_f\left[\expr_r\right]$,
  then $\expr = \ctxt\left[\expr_r\right] \mapsto \ctxt\left[\expr_r'\right]$.
  Otherwise, $\expr = \ctxt\left[\misapp\right]$.

  We assume from now that $\expr_f$ is a value.
  If any of $\sequence{\expr_a}$ is not a value,
  then a similar argument applies.
  Choose $\expr_c$ as the leftmost non-value argument,
  so that $\parens{\sequence{\expr_a}} = \parens{\sequence{\val_a} \; \expr_c \; \sequence{\expr_c'}}$.
  Then induction hypothesis implies that $\expr_c$ is a context $\ctxt_c$
  filled by either the redex $\expr_r$ or the misapplication $\misapp$.
  Then the context $\ctxt =$
  {\tt ($\expr_f$ $\val_a$ $\sequence{}$ $\ctxt_c$ $\expr'_c$ $\sequence{}$)}.
  If $\expr_c = \ctxt_x\left[\expr_r\right] \mapsto \ctxt_x\left[\expr_r'\right]$,
  then $\ctxt\left[\expr_r\right] \mapsto \ctxt\left[\expr_r'\right]$.
  If $\expr_c = \ctxt_x\left[\misapp\right]$, then $\expr = \ctxt\left[\misapp\right]$.

  Having addressed the cases where not all of $\expr_f, \sequence{\expr_a}$ are values,
  we now consider $\expr = \app{\val_f}{\sequence{\val_a}}$.
  By canonical forms, $\val_f$ has the form $\arrlit{\sequence{\func}}{\sequence{\nat_f}}$,
  and uniqueness of typing implies $\ieqv{\idx_f}{\idxshape{\sequence{\nat_f}}}$.
  We proceed by case analysis on $\idx_f, \sequence{\idx_a}$
  (\nb, per the typing derivation, they are all prefix-orderable).
  With one exception---where $\expr$ itself is misapplication of a partial function---%
  each line of argument leads to a particular applicable reduction rule.
  \paragraph{Subcase 1:} $\idx_f = \idxshape{}$, and each $\idx_a = \idxshape{}$.
  Then $\sequence{\nat_f}$ must be the empty sequence,
  and $\val_f$ must contain only a single atom, $\func$.
  So $\val_f = \arrlit{\func}{}$,
  and the respective types of the arguments $\sequence{\val_a}$ are
  $\sequence{\typearray{\type_i}{\idxappend{\idx_a \; \idx_i}}}$,
  or equivalently $\sequence{\typearray{\type_i}{\idx_i}}$.
  If $\func = \primop$,
  then we have the (partially annotated) expression
  $\expr =$
  {\tt ((array () $\primop^{\typefun
      {\sequence{\typearray{\type_i}{\idx_i}}}
      {\typearray{\type_o}{\idx_o}}}$)
    $\val_a^{\typearray{\type_i}{\idx_i}}$
    $\sequence{}$)}.
  If $\primop$ is a partial function
  with $\sequence{\val_a}$ as out-of-domain inputs,
  such as division by zero,
  then the third condition of the progress lemma holds.
  Otherwise, we have a $\delta$ redex.
  On the other hand, if
  $\func = \lam{\sequence{\notevar{\var}{\typearray{\type_i}{\idx_i}}}}{\expr_b}$,
  the annotated expression is
  $\app
  {\arrlit
    {\annotate
      [\typefun
      {\sequence{\typearray{\type_i}{\idx_i}}}
      {\typearray{\type_o}{\idx_o}}]
      {\lam{\sequence{\notevar{\var}{\typearray{\type_i}{\idx_i}}}}{\expr_b}}}{}}
  {\sequence{\annotate[\typearray{\type_i}{\idx_i}]{\val_a}}}$.
  This is a $\beta$ redex.%,
  % which reduces to $\seqsubst{\expr_b}{\var}{\annotate[\typearray{\type_i}{\idx_i}]\val}$.
  \paragraph{Subcase 2:}
  $\idx_f = \idxshape{\sequence{\nat_f}}$,
  and each $\idx_a = \idxshape{\sequence{\nat_f}}$
  for nonempty sequence $\sequence{\nat_f}$.
  Then $\val_f$ is
  $${\annotate
    [\typearray
    {\typefun
      {\sequence{\typearray{\type_i}{\idxshape{\sequence{\nat_i}}}}}
      {\typearray{\type_o}{\idx_o}}}
    {\idxshape{\sequence{\nat_f}}}]
    {\arrlit{\sequence{\func}}{\sequence{\nat_f}}}}$$
  and it is applied to arguments
  $$\sequence{\val_a} =
  \sequence{\annotate
    [\typearray{\type_i}{\idxshape{\sequence{\nat_f} \; \sequence{\nat_i}}}]
    {\arrlit{\sequence{\atval}}{\sequence{\nat_f} \; \sequence{\nat_i}}}}$$
  This is a $\mathit{map}$ redex.%,
  % which reduces to
  % $$\frm
  % {\sequence{
  %     \app
  %     {\arrlit{\func}{}}
  %     {\sequence{\arrlit{\sequence{\atval'}}{\sequence{\nat_i}}}}}}
  % {\sequence{\nat_f}}$$
  % where
  % $\parens{\sequence{\parens{\sequence{\parens{\sequence{\atval'}}}}}}
  % = \mathit{Transpose}\llb\sequence{\mathit{Split}_{\prod\parens{\sequence{\nat_i}}}\llb\sequence{\atval}\rrb}\rrb$
  \paragraph{Subcase 3:}
  $\idx_f, \sequence{\idx_a}$ are not all equal but still prefix orderable.
  The form of $\expr$ is similar to the previous subcase,
  except that each argument array $\val_a$ replaces $\sequence{\nat_f}$
  with its own particular frame shape.
  We then have a $\mathit{lift}$ redex.
\end{sproof}


\begin{lemma}[Preservation]
  \label{Preservation}
  Let $\sEnv,\kEnv,\tEnv$ be a well-formed environment, \ie, $\wfenv{\sEnv}{\kEnv}{\tEnv}$.
  If $\typeof{\sEnv}{\kEnv}{\tEnv}{\expr}{\type}$
  and $\expr \mapsto \expr'$
  then $\typeof{\sEnv}{\kEnv}{\tEnv}{\expr'}{\type}$.
\end{lemma}
\begin{sproof}[We use induction on the derivation of
  $\typeof{\emptyEnv}{\emptyEnv}{\emptyEnv}{\expr}{\type}$.
  An {\tt array} form which can take a reduction step
  must contain a reducible subexpression.
  Many typing rules give rise to subcases where the $\expr$ itself is not a redex
  but contains some subexpression $\expr_r$ which steps to $\expr_r'$.
  In these situations,
  the typing derivation for $\expr_r$ is included in that for $\expr$,
  so replacing that subderivation with one for $\expr_r'$
  (deriving the same type, according to the induction hypothesis)
  produces a derivation of
  $\typeof{\emptyEnv}{\emptyEnv}{\emptyEnv}{\expr'}{\type}$.
  
  The remaining nontrivial subcases each correspond to particular reduction rules.
  As in proving Progress,
  the {\sc T-App} case is split into subcases based on the function and argument frames.
  When frames are non-identical but prefix-compatible,
  the resulting {\it lift} reduction produces an application form with the same principal frame
  and thus the same result type.
  When we have identical non-scalar frames,
  the {\it map} reduction produces a {\tt frame} form
  whose frame shape is equal to the application form's principal frame
  and whose cell shape and atom type is the same as
  the function's return shape and atom type.
  This gives it a type equivalent to that of the {\it map} redex.
  With a scalar principal frame,
  we have a $\delta$ redex (trivial)
  or $\beta$ redex (follows from Lemma \ref{EESub},
  preservation of types under term substitution).
  Reasoning for type and index application forms is similar
  (via Lemma \ref{IESub} and Lemma \ref{TESub} respectively).
  A reducible {\tt unbox} form also
  substitutes a value in for a variable which is intended to have the same type,
  so Lemma \ref{EESub} again ensures the result type is $\type$.]
  We use induction on
  %% the structure of $\expr$.
  the derivation of
  $\typeof{\emptyEnv}{\emptyEnv}{\emptyEnv}{\expr}{\type}$.
  As with Lemma \ref{Progress} (progress),
  we only consider typing rules that can apply to an expression.
  We elide {\sc T-Var} because a variable does not reduce any further.
  \paragraph{Case {\sc Eqv}}
  \[{\infr[T-Eqv]
    {\typeof{\sEnv}{\kEnv}{\tEnv}{\expr}{\type'}
      \qquad \teqv{\type}{\type'}}
    {\typeof{\sEnv}{\kEnv}{\tEnv}{\expr}{\type}}}\]
  By the induction hypothesis,
  $\typeof{\sEnv}{\kEnv}{\tEnv}{\expr'}{\type'}$.
  Then applying {\sc T-Eqv} to that derivation produces
  \[{\infr[T-Eqv]
    {\typeof{\sEnv}{\kEnv}{\tEnv}{\expr'}{\type'}
      \qquad {\teqv{\type}{\type'}}}
    {\typeof{\sEnv}{\kEnv}{\tEnv}{\expr'}{\type}}}\]

  \paragraph{Case {\sc Array}}
  \[{\infr[T-Array]
    {\sequence{\typeof{\sEnv}{\kEnv}{\tEnv}{\atom}{\type_a}}
      \\ \kindof{\sEnv}{\kEnv}{\type_a}{\kindatom}
      \\\\ \mathit{Length}\llb \sequence{\atom} \rrb = \prod{\sequence{\nat}}}
    {\typeof{\sEnv}{\kEnv}{\tEnv}
      {\arrlit{\sequence{\atom}}{\sequence{\nat}}}
      {\typearray{\type_a}{\idxshape{\sequence{\nat}}}}}}\]
  An array literal is not itself a redex,
  so the only way for $\expr$ to reduce to $\expr'$ is
  if some atom $\atom_r$ in the array is reducible.
  Let $\parens{\sequence{\atom_0} \; \atom_r \; \sequence{\atom_1}}
    = \parens{\sequence{\atom}}$.
  The only atom form which may contain an expression
  that takes a reduction step is a box.
  So $\atom_r$ must have the form
  $\dsum
  {\sequence{\idx}}{\expr_s}
  {\typedsum{\sequence{\notevar{\var}{\sort}}}{\type_s}}$,
  where $\expr_s \mapsto \expr_s'$
  and
  $\teqv{\typedsum{\sequence{\notevar{\var}{\sort}}}{\type_s}}{\type_a}$.
  The typing derivation for $\atom_r$ must,
  except perhaps for use of {\sc T-Eqv},
  end with
  \[{\infr[T-Box]
    {\sequence{\sortof{\sEnv}{\idx}{\sort}}
      \qquad
      {\kindof{\sEnv}{\kEnv}{\typedsum{\sequence{\notevar{\var}{\sort}}}{\type_s}}{\kindatom}}
      \\
      \typeof{\sEnv}{\kEnv}{\tEnv}{\expr_s}
      {\seqsubst{\type_s}{\var}{\idx}}}
    {\typeof{\sEnv}{\kEnv}{\tEnv}
      {\dsum
        {\sequence{\idx}}
        {\expr_s}
        {\typedsum{\sequence{\notevar{\var}{\sort}}}{\type_s}}}
      {\typedsum{\sequence{\notevar{\var}{\sort}}}{\type_s}}}}\]
  The induction hypothesis implies that
  $\typeof{\sEnv}{\kEnv}{\tEnv}{\expr_s'}
  {\seqsubst{\type_s}{\var}{\idx}}$,
  so $\atom_r$ can be ascribed the same type
  ${\typedsum{\sequence{\notevar{\var}{\sort}}}{\type_s}}$,
  which is equivalent to $\type_a$.
  Since $\atom_r \mapsto \atom_r'$,
  which still has type $\type_a$,
  we derive
  \[{\infr[T-Array]
    {\hspace{5em}\sequence{\typeof{\sEnv}{\kEnv}{\tEnv}{\atom_0}{\type_a}}
      \\
      \typeof{\sEnv}{\kEnv}{\tEnv}{\atom_r'}{\type_a}
      \\\\
      \sequence{\typeof{\sEnv}{\kEnv}{\tEnv}{\atom_1}{\type_a}}
      \\
      \kindof{\sEnv}{\kEnv}{\type_a}{\kindatom}
      \hspace{5em}
      \\\\
      \mathit{Length}\llb \sequence{\atom} \rrb = \prod{\sequence{\nat}}}
    {\typeof{\sEnv}{\kEnv}{\tEnv}
      {\arrlit
        {\sequence{\atom}}
        {\sequence{\nat}}}
      {\typearray{\type_a}{\idxshape{\sequence{\nat}}}}}}\]
  \paragraph{Case {\sc Frame}}
  \[{\infr[T-Frame]
    {\sequence{\typeof{\sEnv}{\kEnv}{\tEnv}
        {\expr_a}{\typearray{\type_a}{\idx_a}}}
      \\
      \kindof{\sEnv}{\kEnv}{\typearray{\type_a}{\idx_a}}{\kindarray}
      \\
      \mathit{Length}\llb \sequence{\expr} \rrb = \prod{\sequence{\nat_f}}}
    {\typeof{\sEnv}{\kEnv}{\tEnv}
      {\frm{\expr_a}{\sequence{\nat_f}}}
      {\typearray{\type_a}{\idxappend{\idxshape{\sequence{\nat_f}} \; \idx_a}}}}}\]
  As in the {\sc T-Array} case,
  if $\expr$ itself is not a redex,
  then in order for it to reduce to $\expr'$,
  it must contain some reducible cell,
  \ie, there is some $\expr_r \in \sequence{\expr_a}$
  such that $\expr_r \mapsto \expr_r'$.
  Since the type derivation for $\frm{\expr_a}{\sequence{\nat_f}}$
  must ascribe the type $\typearray{\type_a}{\idx_a}$ to $\expr_r$,
  the induction hypothesis gives
  $\typeof{\sEnv}{\kEnv}{\tEnv}{\expr_r'}{\typearray{\type_a}{\idx_a}}$.
  We then patch that result into the type derivation we had.
  With $\sequence{\expr_a} =
  \parens{\sequence{\expr_0} \; \expr_r \; \sequence{\expr_1}}$,
  we derive
  \[{\infr[T-Frame]
    {\sequence{\typeof{\sEnv}{\kEnv}{\tEnv}
        {\expr_0}{\typearray{\type_a}{\idx_a}}}
      \\
      \typeof{\sEnv}{\kEnv}{\tEnv}
             {\expr_r}{\typearray{\type_a}{\idx_a}}
      \\
      \sequence{\typeof{\sEnv}{\kEnv}{\tEnv}
        {\expr_1}{\typearray{\type_a}{\idx_a}}}
      \\
      \kindof{\sEnv}{\kEnv}{\typearray{\type_a}{\idx_a}}{\kindarray}
      \\
      \mathit{Length}\llb \sequence{\expr} \rrb = \prod{\sequence{\nat_f}}}
    {\typeof{\sEnv}{\kEnv}{\tEnv}
      {\frm{\expr_a}{\sequence{\nat_f}}}
      {\typearray{\type_a}{\idxappend{\idxshape{\sequence{\nat_f}} \; \idx_a}}}}}\]
  If $\expr$ \emph{is} a redex,
  it must be a $\mathit{collapse}$ redex.
  Then $\expr$ is a frame of array literals, \ie,
  {\tt (frame
    ($\nat_f \sequence{}$)
    (array ($\nat_c \sequence{}$) $\atval \sequence{}$) $\sequence{}$)},
  and it must step to $\expr' =$
{\tt (array ($\nat_f \sequence{} \nat_c \sequence{}$)
$\mathit{Concat}\llbracket\sequence{(\sequence{\atval})}\rrbracket$)}.
  Each of $\expr$'s cells is a well-typed array literal,
  with type derivation:
  \[{\infr[T-Array]
    {\sequence{\typeof{\sEnv}{\kEnv}{\tEnv}
        {\atval}{\type_a}}
      \qquad
      \kindof{\sEnv}{\kEnv}{\type_a}{\kindatom}
      \\\\
      \mathit{Length}\llb\sequence{\atval}\rrb = \prod{\sequence{\nat_c}}}
    {\typeof{\sEnv}{\kEnv}{\tEnv}
      {\arrlit{\sequence{\atval}}{\sequence{\nat_c}}}
      {\typearray{\type_a}{\idxshape{\sequence{\nat_c}}}}}}\]
  The typing derivation for each cell in the {\tt frame}
  requires that $\mathit{Length}\llb\sequence{\atval}\rrb = \prod{\sequence{\nat_c}}$.
  Concatenating $\prod{\sequence{\nat_f}}$ sequences of atomic values
  each of which contains $\prod{\sequence{\nat_c}}$ elements
  gives a sequence whose length is $\prod{\parens{\sequence{\nat_f}\sequence{\nat_c}}}$.
  So we can derive a similar type for the collapsed array:
  \[{\infr[T-Array]
    {\sequence{\typeof{\sEnv}{\kEnv}{\tEnv}
        {\atval}{\type_a}}
      \\
      \kindof{\sEnv}{\kEnv}{\type_a}{\kindatom}
      \\\\
      \mathit{Length}\llb\mathit{Concat}\llb\sequence{\parens{\sequence{\atval}}}\rrb\rrb
      = \prod{\parens{\sequence{\nat_f}\sequence{\nat_c}}}}
    {\parbox{0.65\textwidth}
      {\centering ${\sEnv};{\kEnv};{\tEnv}\vdash
        {\arrlit
          {\mathit{Concat}\llb\sequence{\parens{\sequence{\atval}}}\rrb}
          {\sequence{\nat_f}\sequence{\nat_c}}}$
        $:{\typearray{\type_a}{\idxshape{\sequence{\nat_f}\sequence{\nat_c}}}}$}}}\]

  Strictly speaking, our goal is to ascribe the type
  {\tt (Arr $\type_a$ (++ (Shp $\nat_f \sequence{}$) $\idx_a$))}  $=$
  {\tt (Arr $\type_a$ (++ (Shp $\nat_f \sequence{}$) (Shp $\nat_c \sequence{}$)))}
  rather than the type we have just derived,
  but they are equivalent via {\sc TEqv-Array} because
  {\tt (++ (Shp $\nat_f \sequence{}$) (Shp $\nat_c \sequence{}$))
    $\theq$
    (Shp $\nat_f \sequence{} \nat_c \sequence{}$)}.
  So {\sc T-Eqv} completes the derivation for
${\sEnv};{\kEnv};{\tEnv}$
$\vdash$
$\expr'$
$:$
{\tt (Arr $\type_a$ (++ (Shp $\nat_f \sequence{}$) $\idx_a$))}.

  \paragraph{Case {\sc TApp}}
  \[{\infr[T-TApp]
    {\typeof{\sEnv}{\kEnv}{\tEnv}
      {\expr_f}
      {\typearray
        {\typeuniv
          {\sequence{\notevar{\var}{\kind}}}
          {\typearray{\type_u}{\idx_u}}}
        {\idx_f}}
      \qquad
      \sequence{\kindof{\sEnv}{\kEnv}{\type_a}{\kind}}}
    {\typeof{\sEnv}{\kEnv}{\tEnv}
      {\tapp{\expr_f}{\sequence{\type_a}}}
      {\typearray
        {\seqsubst{\type_u}{\var}{\type_a}}
        {\idxappend{\idx_f \; \idx_u}}}}}\]
  As in previous cases, if $\expr_f \mapsto \expr_f'$,
  the induction hypothesis gives us a derivation
  $$\typeof{\sEnv}{\kEnv}{\tEnv}
  {\expr_f'}
  {\typearray
    {\typeuniv
      {\sequence{\notevar{\var}{\kind}}}
      {\typearray{\type_u}{\idx_u}}}
    {\idx_f}}$$
  which we can then use for
  {\tt (t-app $\expr_f'$ $\type_a \sequence{}$)}.

  Otherwise, $\expr \mapsto \expr'$ is only possible
  if we have a $\mathit{t\beta}$ redex,
  with
  $\expr_f =$ {\tt (array ($\nat_f \sequence{}$) (T$\ttlm{\sequence{\ttparens{\var_u\; \kind}}}$ $\expr_u$) $\sequence{}$)}
  and $\idx_f = \idxshape{\sequence{\nat_f}}$.
  Following $\mathit{t\beta}$ reduction, $\expr' =$
  {\tt (frame ($\nat_f \sequence{}$) $\seqsubst{\expr_u}{\var_u}{\type_a} \sequence{}$)}.
  Since $\expr_f$ is a well-typed array of type abstractions,
  it must be the case that
  $\typeof{\sEnv}{\kEnv, \sequence{\haskind{\var_u}{\kind}}}{\tEnv}
  {\expr_u}
  {\typearray{\type_u}{\idx_u}}$
  for each of the abstraction bodies, $\sequence{\expr_u}$.
  Lemma \ref{TESub} (preservation of types under type substitution)
  implies that we can give $\seqsubst{\expr_u}{\var_u}{\type_a}$
  the ``same'' type as $\expr_u$,
  that is,
  $$\typeof{\sEnv}{\kEnv}{\seqsubst{\tEnv}{\var_u}{\type_a}}
  {\seqsubst{\expr_u}{\var_u}{\type_a}}
  {\typearray{\seqsubst{\type_u}{\var_u}{\type_a}}{\idx_u}}$$

  The type derivation for $\expr_f$ requires that
  the individual type abstractions' types all be
  equivalent to $\typeuniv{\sequence{\notevar{\var}{\kind}}}{\type_u}$
  despite possibly being written to bind different type names,
  so each resulting $\seqsubst{\type_u}{\var_u}{\type_a}$
  is equivalent to $\seqsubst{\type_u}{\var}{\type_a}$.
  Per Barendregt's convention,
  an abstraction's type variables $\sequence{\var_u}$
  are not used in the range of $\tEnv$,
  the type environment used in checking that abstraction
  (otherwise, the environment structure $\sEnv; \kEnv; \tEnv$
  would be ill-formed),
  so $\seqsubst{\tEnv}{\var_u}{\type_a} = \tEnv$.
  Theorem \ref{wellkinded} (ascription of well-kinded types)
  implies that the pieces used to form the result type,
  $\seqsubst{\type_u}{\var}{\type_a}$ and $\idx_u$,
  are well-formed at kind $\kindatom$ and sort $\sortshp$ respectively,
  so the cell type $\typearray{\seqsubst{\type_u}{\var}{\type_a}}{\idx_u}$
  has kind $\kindarray$.
  Finally, the number of result cells matches
  the number of atoms in $\expr_f$,
  which is the product of $\expr_f$'s dimensions $\sequence{\nat_f}$.
  Thus we can derive
  \[{\infr[T-Frame]
    {\sequence{
        \typeof{\sEnv}{\kEnv}{\tEnv}
        {\seqsubst{\expr_u}{\var_u}{\type_a}}
        {\typearray
          {\seqsubst{\type_u}{\var}{\type_a}}
          {\idx_u}}}
      \\\\
      \kindof
          {\sEnv}{\kEnv}
          {\typearray
            {\seqsubst{\type_u}{\var}{\type_a}}
            {\idx_u}}
          {\kindarray}
      \\\\
      \mathit{Length}\llb \sequence{\expr_u} \rrb = \prod{\sequence{\nat}}}
    {\parbox{0.6\textwidth}
      {\centering
        ${\sEnv};{\kEnv};{\tEnv}\vdash
        {\frm
          {\sequence{\seqsubst{\expr_u}{\var_u}{\type_a}}}
          {\sequence{\nat_f}}}$
        $: {\typearray
          {\seqsubst{\type_u}{\var}{\type_a}}
          {\idxappend{\idxshape{\sequence{\nat}} \; \idx_u}}}
        $}}}\]

  \paragraph{Case {\sc IApp}}
  \newcommand{\ibetasub}[1]{\seqsubst{#1}{\var}{\idx_a}}
  \newcommand{\ibetasubp}[1]{\seqsubst{#1}{\var'}{\idx_a}}
  \[\infr[T-IApp]
  {\typeof{\sEnv}{\kEnv}{\tEnv}
    {\expr_f}
    {\typearray
      {\typedprod
        {\sequence{\notevar{\var}{\sort}}}
        {\typearray{\type_p}{\idx_p}}}
      {\idx_f}}}
  {\typeof{\sEnv}{\kEnv}{\tEnv}
    {\iapp{\expr_f}{\idx_a}}
    {\typearray{\ibetasub{\type_p}}{\idxappend{\idx_f \; \ibetasub{\idx_p}}}}}\]
  If $\expr_f \mapsto \expr_f'$,
  then the induction hypothesis implies that
  ${\sEnv};{\kEnv};{\tEnv}\vdash$
  ${\expr_f'}$
  {\tt (Arr (Pi (($\var$ $\sort$) $\sequence{}$) (Arr $\type_p$ $\idx_p$)) $\idx_f$)}.
  Then the type derivation for $\expr_f'$ can be substituted into the original derivation for $\expr$,
  ascribing the same type
  to the reduced term
  ${\iapp{\expr_f'}{\sequence{\idx_a}}}$.

  Otherwise, $\expr$ must be an $\mathit{i\beta}$ redex,
  with $\expr_f = \arrlit{\sequence{\atval}}{\sequence{\nat_f}}$,
  each $\atval$ of the form $\ilam{\sequence{\notevar{\var'}{\sort}}}{\val}$,
  and $\ieqv{\idx_f}{\idxshape{\sequence{\nat_f}}}$.
  We must also be able to derive, for each $\atval$,
  $\typeof{\sEnv}{\kEnv}{\tEnv}
  {\atval}{\typedprod{\sequence{\notevar{\var}{\sort}}}{\typearray{\type_p}{\idx_p}}}$.
  Note that $\sequence{\var'}$ might not be the same variables as $\sequence{\var}$.
  If not, then the type derivation for $\atval$ must end with {\sc T-Eqv}
  relating the required array type $\typearray{\type_p}{\idx_p}$
  and an actually derived array type $\typearray{\type_p'}{\idx_p'}$.
  That is, use of {\sc T-Eqv} requires deriving both
  \[{\infr[T-ILam]
    {\typeof{\sEnv, \sequence{\hassort{\var'}{\sort}}}{\kEnv}{\tEnv}
      {\val}
      {\typearray{\type_p'}{\idx_p'}}}
    {\typeof{\sEnv}{\kEnv}{\tEnv}
      {\ilam{\sequence{\notevar{\var'}{\sort}}}{\val}}
      {\typedprod
        {\sequence{\notevar{\var'}{\sort}}}
        {\typearray{\type_p'}{\idx_p'}}}}}\]
  and, with free variables $\sequence{\var_f}$,
  \[{\infr[TEqv-Pi]
    {\infr[TEqv-Array]
      {\teqv
        {\seqsubst{\type_p}{\var}{\var_f}}
        {\seqsubst{\type_p'}{\var'}{\var_f}}
        \\\\
        \ieqv{\seqsubst{\idx_p}{\var}{\var_f}}{\seqsubst{\idx_p'}{\var'}{\var_f}}}
      {\teqv
        {\seqsubst{\typearray{\type_p}{\idx_p}}{\var}{\var_f}}
        {\seqsubst{\typearray{\type_p'}{\idx_p'}}{\var'}{\var_f}}}}
    {\teqv
      {\typedprod{\sequence{\notevar{\var}{\sort}}}{\typearray{\type_p}{\idx_p}}}
      {\typedprod{\sequence{\notevar{\var'}{\sort}}}{\typearray{\type_p'}{\idx_p'}}}}}\]

  We now ascribe a type to $\ibetasubp{\val}$.
  Previously, $\val$ was given the type $\typearray{\type_p'}{\idx_p'}$,
  so Lemma \ref{IESub} (preservation of types under index substitution) implies
  $\typeof{\sEnv}{\kEnv}{\ibetasubp{\tEnv}}
  {\ibetasubp{\val}}{\ibetasubp{\typearray{\type_p'}{\idx_p'}}}$.
  Well-formedness of the environment% $\sEnv;\kEnv;\tEnv$
  means that no free index variables are used in the range of $\tEnv$,
  so $\ibetasubp{\tEnv} = \tEnv$.
  This simplifies our typing result to
  $\typeof{\sEnv}{\kEnv}{\tEnv}
  {\ibetasubp{\val}}{\ibetasubp{\typearray{\type_p'}{\idx_p'}}}$.
  The ascribed type is equal to $\typearray{\ibetasubp{\type_p'}}{\ibetasubp{\idx_p'}}$.
  We must show that it is equivalent to
  $\typearray{\ibetasub{\type_p}}{\ibetasub{\idx_p}}$.

  We consider
  $\seqsubst{\seqsubst{\type_p}{\var}{\var_f}}{\var_f}{\idx_a}$
  and
  $\seqsubst{\seqsubst{\type_p'}{\var'}{\var_f}}{\var_f}{\idx_a}$,
  which can be simplified to
  $\seqsubst{\type_p}{\var}{\idx_a}$
  and
  $\seqsubst{\type_p'}{\var'}{\idx_a}$
  respectively.
  Since index substitution preserves type equivalence (per Lemma \ref{EqlIdxEqvType}),
  our earlier $\teqv{\seqsubst{\type_p}{\var}{\var_f}}{\seqsubst{\type_p'}{\var'}{\var_f}}$
  implies that $\teqv{\seqsubst{\type_p}{\var}{\idx_a}}{\seqsubst{\type_p'}{\var'}{\idx_a}}$.
  Similarly, preservation of index equality under substitution implies
  $\ieqv{\seqsubst{\idx_p}{\var}{\idx_a}}{\seqsubst{\idx_p'}{\var'}{\idx_a}}$.
  So we can derive the type equivalence
  \[{\infr[TEqv-Array]
    {\teqv
      {\seqsubst{\type_p}{\var}{\idx_a}}
      {\seqsubst{\type_p'}{\var'}{\idx_a}}
      \\\\
      \ieqv{\seqsubst{\idx_p}{\var}{\idx_a}}{\seqsubst{\idx_p'}{\var'}{\idx_a}}}
    {\teqv
      {\ibetasub{\typearray{\type_p}{\idx_p}}}
      {\ibetasubp{\typearray{\type_p'}{\idx_p'}}}}}\]
  which then allows us to derive for each of $\sequence{\val}$
  \[{\infr[T-Eqv]
    {\typeof{\sEnv}{\kEnv}{\tEnv}
      {\ibetasubp{\val}}
      {\ibetasub{\typearray{\type_p'}{\idx_p'}}}
      \\\\
      \teqv
      {\ibetasubp{\typearray{\type_p'}{\idx_p'}}}
      {\ibetasub{\typearray{\type_p}{\idx_p}}}}
    {\typeof{\sEnv}{\kEnv}{\tEnv}
      {\ibetasubp{\val}}
      {\ibetasub{\typearray{\type_p}{\idx_p}}}}}\]
  This enables a type derivation for
  $\expr' = \frm{\sequence{\ibetasubp{\val}}}{\sequence{\nat_f}}$:
  \[{\infr[T-Frame]
    {\sequence{\typeof{\sEnv}{\kEnv}{\tEnv}
        {\ibetasubp{\val}}
        {\ibetasub{\typearray{\type_p}{\idx_p}}}}
      \\\\
      \mathit{Length}\llb \sequence{\val} \rrb = \prod{\parens{\sequence{\nat_f}}}}
    {\typeof{\sEnv}{\kEnv}{\tEnv}
      {\expr'}
      {\typearray
        {\ibetasub{\type_p}}
        {\idxappend{\idxshape{\sequence{\nat_f}} \; \ibetasub{\idx_p}}}}}}\]

  \paragraph{Case {\sc Unbox}}
  \[{\infr[T-Unbox]
    {\typeof{\sEnv}{\kEnv}{\tEnv}{\expr_s}
      {\typearray{\typedsum{\sequence{\notevar{\var_i'}{\sort}}}{\type_s}}{\idx_s}}
      \\\\
      \typeof{\sEnv, \sequence{\hassort{\var_i}{\sort}}}
      {\kEnv}
      {\tEnv, \hastype{\var_e}{\seqsubst{\type_s}{\var_i'}{\var_i}}}
      {\expr_b}{\typearray{\type_b}{\idx_b}}
      \\\\
      \kindof{\sEnv}{\kEnv}{\typearray{\type_b}{\idx_b}}{\kindarray}}
    {\typeof{\sEnv}{\kEnv}{\tEnv}
      {\dproj{\sequence{\var_i}}{\var_e}{\expr_s}{\expr_b}}
      {\typearray{\type_b}{\idxappend{\idx_s \; \idx_b}}}}}\]
  If $\expr$ is reducible because $\expr_s \mapsto \expr_s'$,
  then the induction hypothesis implies
  $\typeof{\sEnv}{\kEnv}{\tEnv}{\expr_s'}
  {\typearray{\typedsum{\sequence{\notevar{\var_i'}{\sort}}}{\type_s}}{\idx_s}}$.
  This can be used to adapt the original type derivation for $\expr$
  to fit $\expr' = {\dproj{\sequence{\var_i}}{\var_e}{\expr_s'}{\expr_b}}$:
  \[{\infr[T-Unbox]
    {\typeof{\sEnv}{\kEnv}{\tEnv}{\expr_s'}
      {\typearray{\typedsum{\sequence{\notevar{\var_i'}{\sort}}}{\type_s}}{\idx_s}}
      \\\\
      \typeof{\sEnv, \sequence{\hassort{\var_i}{\sort}}}
      {\kEnv}
      {\tEnv, \hastype{\var_e}{\seqsubst{\type_s}{\var_i'}{\var_i}}}
      {\expr_b}{\typearray{\type_b}{\idx_b}}
      \\\\
      \kindof{\sEnv}{\kEnv}{\typearray{\type_b}{\idx_b}}{\kindarray}}
    {\typeof{\sEnv}{\kEnv}{\tEnv}
      {\dproj{\sequence{\var_i}}{\var_e}{\expr_s'}{\expr_b}}
      {\typearray{\type_b}{\idxappend{\idx_s \; \idx_b}}}}}\]
  Otherwise, $\expr$ must be an $\mathit{unbox}$ redex,
  where all of the following hold:
  \begin{enumerate}
  \item{$\expr_s$ is of the form
      $\arrlit{\sequence{\dsum{\sequence{\idx_s}}{\val_s}{\type_\sigma}}}{\sequence{\nat_s}}$}
  \item{$\type_\sigma =
      \teqv
      {\typedsum{\sequence{\notevar{\var_i''}{\sort}}}{\type_s'}}
      {\typedsum{\sequence{\notevar{\var_i'}{\sort}}}{\type_s}}$}
  \item{$\idx_s = \idxshape{\sequence{\nat_s}}$}
  \item{$\mathit{Length}\llb \sequence{\dsum{\sequence{\idx_s}}{\val_s}{\type_\sigma}} \rrb
      = \prod{\sequence{\nat_s}}$}
  \end{enumerate}

  The type derivation for $\expr_s$ must include
  ascription of $\type_\sigma$ to each box within the array:
  \[{\infr[T-Box]
    {\sequence{\sortof{\sEnv}{\idx_\sigma}{\sort}}
      \qquad
      {\kindof{\sEnv}{\kEnv}{\type_\sigma}{\kindatom}}
      \\
      \typeof{\sEnv}{\kEnv}{\tEnv}{\val_\sigma}{\seqsubst{\type_s'}{\var_i''}{\idx_\sigma}}}
    {\typeof{\sEnv}{\kEnv}{\tEnv}
      {\dsum{\sequence{\idx_\sigma}}{\val_\sigma}{\type_\sigma}}
      {\type_\sigma}}}\]

  Equivalence of $\type_\sigma$ and $\typedsum{\sequence{\notevar{\var_i'}{\sort}}}{\type_s}$
  means that each box's $\val_\sigma$ can also be typed as $\seqsubst{\type_s}{\var_i'}{\idx_\sigma}$.

  The resulting term $\expr'$ is
  $\multisubst{\expr_b}{\sequence{\singlesubst{\var_i}{\idx_\sigma}},\singlesubst{\var_e}{\val_\sigma}}
  = \subst{\seqsubst{\expr_b}{\var_i}{\idx_\sigma}}{\var_e}{\val_\sigma}$.
  Since we have a type derivation for $\expr_b$ in an extended environment,
  we will apply Lemmas \ref{IESub} and \ref{EESub}
  (preservation of types under index and expression substitution)
  to produce a type derivation in the original environment.

  Using the fact that each of $\sequence{\idx_\sigma}$ has its required sort
  (necessary for the previous derivation),
  Lemma \ref{IESub} gives us
  $\typeof{\sEnv}{\kEnv}{\seqsubst{\parens{\tEnv, \hastype{\var_e}{\seqsubst{\type_s}{\var_i'}{\var_i}}}}{\var_i}{\idx_\sigma}}
  {\seqsubst{\expr_b}{\var_i}{\idx_\sigma}}
  {\seqsubst{\parens{\seqsubst{\type_b}{\var_i}{\var_i'}}}{\var_i}{\idx_\sigma}}$.
  Well-formedness of the original environment means $\tEnv$ does not refer to the
  new index variables $\sequence{\var_i}$,
  which are instead bound by the {\tt unbox} expression.
  So the substituted environment
  $\seqsubst{\parens{\tEnv, \hastype{\var_e}{\seqsubst{\type_s}{\var_i'}{\var_i}}}}{\var_i}{\idx_\sigma}$
  is equal to $\tEnv, \hastype{\var_e}{\seqsubst{\type_s}{\var_i'}{\idx_\sigma}}$,
  turning the type derivation into
  $\typeof{\sEnv}{\kEnv}{\tEnv, \hastype{\var_e}{\seqsubst{\type_s}{\var_i'}{\idx_\sigma}}}
  {\seqsubst{\expr_b}{\var_i}{\idx_\sigma}}{\seqsubst{\type_b}{\var_i}{\var_i'}}$.
  The index variables $\sequence{\var_i}$ are not bound in $\sEnv$,
  but we still have $\kindof{\sEnv}{\kEnv}{\typearray{\type_b}{\idx_b}}{\kindarray}$.
  This means $\type_b$ and $\idx_b$ must be well-formed without those bindings,
  \ie, none of $\sequence{\var_i}$ appear free in $\type_b$ or $\idx_b$.
  Thus $\seqsubst{\typearray{\type_b}{\idx_b}}{\var_i}{\idx_s}$ is simply $\typearray{\type_b}{\idx_b}$.
  Our derivation now concludes
  $\typeof{\sEnv}{\kEnv}{\tEnv, \hastype{\var_e}{\seqsubst{\type_s}{\var_i'}{\idx_\sigma}}}
  {\seqsubst{\expr_b}{\var_i}{\idx_\sigma}}{\typearray{\type_b}{\idx_b}}$.

  Since $\typeof{\sEnv}{\kEnv}{\tEnv}{\val_s}{\seqsubst{\type_s}{\var_i'}{\idx_\sigma}}$,
  Lemma \ref{EESub} produces a derivation of
  $\typeof{\sEnv}{\kEnv}{\tEnv}
  {\subst{\seqsubst{\expr_b}{\var_i}{\idx_\sigma}}{\var_e}{\val_\sigma}}
  {\typearray{\type_b}{\idx_b}}$
  for each box's index contents $\sequence{\idx_\sigma}$ and term contents $\val_\sigma$.
  We can therefore type the result of the $\mathit{unbox}$ reduction step,
  {\tt (frame ($\nat_s$ $\sequence{}$)
    $\subst{\seqsubst{\expr_b}{\var_i}{\idx_\sigma}}{\var_e}{\val_\sigma}$ $\sequence{}$)}
  as follows:
  \[{\infr[T-Frame]
    {\sequence{\typeof{\sEnv}{\kEnv}{\tEnv}
        {\subst{\seqsubst{\expr_b}{\var_i}{\idx_\sigma}}{\var_e}{\val_\sigma}}
        {\typearray{\type_b}{\idx_b}}}
      \\\\
      \kindof{\sEnv}{\kEnv}{\typearray{\type_b}{\idx_b}}{\kindarray}
    \\\\
    \mathit{Length}\llb{\subst{\seqsubst{\expr_b}{\var_i}{\idx_\sigma}}{\var_e}{\val_\sigma}}\rrb
    = \prod{\sequence{\nat_s}}}
  {\parbox{0.6\textwidth}
    {\centering
      \({\sEnv};{\kEnv};{\tEnv}\vdash\)
      \({\frm{\subst{\seqsubst{\expr_b}{\var_i}{\idx_\sigma}}{\var_e}{\val_\sigma}}{\sequence{\nat_s}}}\)
      \({: \typearray{\type_b}{\idxappend{\idx_s}{\idx_b}}}\)}}
  }\]

  \paragraph{Case {\sc App}}
  \[{\infr[T-App]
    {\typeof{\sEnv}{\kEnv}{\tEnv}
      {\expr_f}
      {\typearray
        {\typefun
          {\sequence{\typearray{\type_i}{\idx_i}}}
          {\typearray{\type_o}{\idx_o}}}
        {\idx_f}}
      \\\\
      \sequence
      {\typeof{\sEnv}{\kEnv}{\tEnv}
        {\expr_a}
        {\typearray{\type_i}{\idxappend{\idx_a \; \idx_i}}}}
      \qquad
      \idx_p = \mathit{Max}\llb \idx_f \; \sequence{\idx_a} \rrb}
    {\typeof{\sEnv}{\kEnv}{\tEnv}
      {\app{\expr_f}{\sequence{\expr_a}}}
      {\typearray{\type_o}{\idxappend{\idx_p \; \idx_o}}}}}\]
  If either $\expr_f \mapsto \expr_f'$
  or there is some argument $\expr_A \in \sequence{\expr_a}$
  such that $\expr_A \mapsto \expr_A'$,
  then the induction hypothesis implies that
  the result $\expr_f'$ or $\expr_A'$ retains the same type as $\expr_f$ or $\expr_A$.
  The type derivation for $\expr$ can be updated
  with the new sub-derivation for $\expr_f$ or $\expr_A$
  to derive the same type for $\expr'$.

  Otherwise, $\expr$ must itself be a redex.
  The possible reductions for an application form
  are $\mathit{lift}$, $\mathit{map}$, $\beta$, and $\delta$.
  \paragraph{Subcase 1:}
  If $\expr$ is a $\mathit{lift}$ redex,
  then it has the form
  {\tt
    ((array ($\nat_f$ $\sequence{}$) $\atval_f$ $\sequence{}$)
    (array ($\nat_a$ $\sequence{}$ $\nat_i$ $\sequence{}$)
    $\atval_a$ $\sequence{}$) $\sequence{}$)}.
  In order to match the left hand side of the $\mathit{lift}$ rule,
  it must be the case that that
  $\ieqv{\idx_f}{\idxshape{\sequence{\nat_f}}}$
  and that for each argument's own type derivation,
  $\ieqv{\idx_a}{\idxshape{\sequence{\nat_a}}}$.
  The frame portion of each array's shape
  (\ie, $\sequence{\nat_f}$ for function position
  and $\sequence{\nat_a}$ for argument position)
  is replaced with the principal frame, $\sequence{\nat_p}$.
  The individual atoms used in the new function and arugment arrays
  all come from their corresponding original arrays
  and therefore retain the same types.
  That is, the atoms $\sequence{\atval_f'}$ in the new function array
  $\expr_f' = \arrlit
  {\mathit{Concat}\llb
    \mathit{Rep}_{\nat_{\mathit{fe}}}\llb
    \mathit{Split}_{1}\llb \sequence{\atval_f}
    \rrb\rrb\rrb}
  {\sequence{\nat_p}}$
  all have type
  $\typefun
  {\sequence{\typearray{\type_i}{\idx_i}}}
  {\typearray{\type_o}{\idx_o}}$,
  and the atoms $\sequence{\atval_a'}$ in any new argument array
  $\expr_a' = \arrlit
  {\mathit{Concat}\llb
    \mathit{Rep}_{\nat_{\mathit{ae}}}\llb
    \mathit{Split}_{\nat_{\mathit{ac}}}\llb \sequence{\atval_a}
    \rrb\rrb\rrb}
  {\sequence{\nat_p} \; \sequence{\nat_i}}$
  are all typed as the corresponding $\type_i$.
  While $\mathit{Split}$ breaks a sequence into subsequences
  (preserving the total number of elements)
  and $\mathit{Concat}$ merges a sequence-of-sequences into a single sequence
  (also preserving the total number of elements),
  $\mathit{Rep}_{\nat}$ produces $\nat$ copies of each element.
  The factor used by $\mathit{lift}$ reduction for replication of each array's cells is the quotient of
  the number of cells in the principal frame and the number of cells in the original array.
  Since the original frame shape is a prefix of the principal frame,
  divisibility is guaranteed.
  This also ensures that the number of atoms in each new argument array is
  $\nat_{ae}*\parens{\prod\sequence{\nat_a}}*\parens{\prod\sequence{\nat_i}}
  = \parens{\prod\sequence{\nat_p}}*\parens{\prod\sequence{\nat_i}}
  = \prod{\parens{\sequence{\nat_p}\;\sequence{\nat_i}}}$.
  Similarly, the number of atoms in the new function array is
  $\nat_{fe}*\parens{\prod\sequence{\nat_f}} = \prod{\sequence{\nat_p}}$.
  Based on the types and quantity of the function and argument arrays' atoms,
  we can derive
  \[{\infr[T-App]
    {\typeof{\sEnv}{\kEnv}{\tEnv}
      {\expr_f'}
      {\typearray
        {\typefun
          {\sequence{\typearray{\type_i}{\idx_i}}}
          {\typearray{\type_o}{\idx_o}}}
        {\idxshape{\sequence{\nat_p}}}}
      \\
      \sequence{\typeof{\sEnv}{\kEnv}{\tEnv}
        {\expr_a'}
        {\typearray{\type_i}{\idxappend{\idxshape{\sequence{\nat_p}} \; \idx_i}}}}
      \\\\
    \idx_p = \mathit{Max} \llb \idx_p \; \sequence{\idx_p} \rrb}
    {\typeof{\sEnv}{\kEnv}{\tEnv}
      {\expr'}
      {\typearray{\type_o}{\idxappend{\idx_p \; \idx_o}}}}}\]
  \paragraph{Subcase 2:}
  If $\expr$ is a $\mathit{map}$ redex,
  then it has the form
  {\tt
    ((array ($\sequence{\nat_f}$) $\sequence{\atom_f}$)
    (array ($\sequence{\nat_f}$ $\sequence{\nat_a}$) $\sequence{\atom_a}$) $\sequence{}$)}.
  % $\app
  % {\arrlit{\sequence{\atom_f}}{\sequence{\nat_f}}}
  % {\sequence{\arrlit{\sequence{\atom_a}}{\sequence{\nat_f} \; \sequence{\nat_a}}}}$.
  Matching the left hand side of the $\mathit{map}$ rule requires that
  $\ieqv{\idx_i}{\idxshape{\sequence{\nat_i}}}$.
  Each argument $\expr_a = \arrlit{\sequence{\atom_a}}{\sequence{\nat_f} \; \sequence{\nat_a}}$
  has its sequence of atoms split into segments whose length is
  the product of $\sequence{\nat_i}$, the corresponding input type's dimensions.
  Transposition groups the first segment from each argument,
  then the second segment from each argument, and so on.
  So the nested sequence
  $\sequence{\parens{\sequence{\parens{\sequence{\atval_c}}}}}$
  has for each $j^{\text{th}}$ $\parens{\sequence{\parens{\sequence{\atval_c}}}}$
  a sequence of atoms corresponding to each of the function input types:
  the $(j,k)^{\text{th}}$ atom sequence contains the atoms which make up
  the $j^{\text{th}}$ cell of the $k^{\text{th}}$ argument.
  Since the length of this single $\sequence{\atval_c}$ is
  the product of the corresponding input type's dimensions
  and they all have the same type as the atoms in the corresponding input type,
  the array literal built from them,
  $\arrlit{\sequence{\atval_c}}{\sequence{\nat_i}}$
  has type $\typearray{\type_i}{\idxshape{\sequence{\nat_i}}}$.
  Each function array $\arrlit{\atval_f}{}$
  contains a single atom taken from the original $\expr_f$,
  implying that it must have type
  $\typefun
  {\sequence{\typearray{\type_i}{\idx_i}}}
  {\typearray{\type_o}{\idx_o}}$,
  which is equivalent to
  $\typefun
  {\sequence{\typearray{\type_i}{\idxshape{\sequence{\nat_i}}}}}
  {\typearray{\type_o}{\idx_o}}$.
  Since the single-cell argument arrays' types all match
  the singleton function array's input types
  with principal frame of $\idxshape{}$,
  each application form itself has type $\typearray{\type_o}{\idx_o}$.
  That is, each result-cell application form
  $\expr_c =$
  {\tt
    ((array () $\atval_f$)
    (array ($\sequence{\nat_i}$) $\sequence{\atval_c}$) $\sequence{}$)}
  in the resulting {\tt frame} form
  can be typed as:
  \[\infr[T-App]
  {\parbox{0.6\textwidth}
    {\centering
      \({\sEnv};{\kEnv};{\tEnv}\vdash\)
      ${\arrlit{\atval_f}{}}$\\
      $:{\typearray
        {\typefun
          {\sequence{\typearray{\type_i}{\idx_i}}}
          {\typearray{\type_o}{\idx_o}}}
        {\idxshape{}}}$}
    \\\\\\
    \sequence{\typeof{\sEnv}{\kEnv}{\tEnv}
      {\arrlit{\sequence{\atval_c}}{\sequence{\nat_i}}}
      {\typearray{\type_i}{\idx_i}}}
    \\\\\\
    \idx_p = \idxshape{}}
  {\typeof{\sEnv}{\kEnv}{\tEnv}
    {\expr_c}
    {\typearray{\type_o}{\idx_o}}}\]
  Recall that our goal type in this case is
  $\type =$
  $\typearray{\type_o}{\idxappend{\idx_p\;\idx_o}}$,
  where $\idx_p = \idxshape{\sequence{\nat_f}}$.
  That is, $\type$ is
  $\typearray{\type_o}{\idxappend{\idxshape{\sequence{\nat_f}}\;\idx_o}}$.
  Using the above derivations for the result cells $\sequence{\expr_c}$, we then derive
  \[{\infr[T-Frame]
    {{\sequence{\typeof{\sEnv}{\kEnv}{\tEnv}
          {\expr_c}
          {\typearray{\type_o}{\idx_o}}}}
      \\
      \kindof{\sEnv}{\kEnv}{\type}{\kindarray}}
    {\typeof{\sEnv}{\kEnv}{\tEnv}
      {\frm
        {\sequence{\expr_c}}
        {\sequence{\nat_f}}}
      {\typearray{\type_o}{\idxappend{\idxshape{\sequence{\nat_f}}\;\idx_o}}}}}\]
  \paragraph{Subcase 3:}
  If $\expr$ is a $\beta$ redex,
  then $\expr_f$ must have the form
  {\tt
    (array () (\cl (($\var$ (Arr $\type_i$ $\idx_i$)) $\sequence{}$) $\expr_0$))}.
  Then $\expr \mapsto_\beta \expr' = \seqsubst{\expr_o}{\var}{\expr_a}$.
  We also know for each of $\sequence{\expr_a}$ that
  $\typeof{\sEnv}{\kEnv}{\tEnv}{\expr_a}{\typearray{\type_i}{\idx_i}}$.
  In the type derivation for $\expr$ itself,
  we must have $\idx_f$ and every $\idx_a$ equal to $\idxshape{}$,
  so $\idx_p = \idxshape{}$.
  Thus
  $\teqv
  {\typearray{\type_o}{\idxappend{\idx_p \; \idx_o}}}
  {\typearray{\type_o}{\idx_o}}$
  The type of the function atom in $\expr_f$ is derived either by {\sc T-Lam},
  in which case, the derivation must ascribe $\typearray{\type_o}{\idx_o}$ to $\expr_o$,
  or by {\sc T-Eqv},
  in which case there must be some earlier use of {\sc T-Lam}
  which ascribes some type equivalent type to $\expr_o$.
  Either way, we have
  $\typeof{\sEnv}{\kEnv}{\tEnv, \sequence{\hastype{\var}{\typearray{\type_i}{\idx_i}}}}
  {\expr_o}
  {\typearray{\type_o}{\idx_o}}$.
  The substitution lemma (Lemma \ref{EESub}) then gives
  $\typeof{\sEnv}{\kEnv}{\tEnv}{\expr'}{\typearray{\type_o}{\idx_o}}$.
  \paragraph{Subcase 4:}
  If $\expr$ is a $\delta$ redex,
  then the required result follows from
  the requirements for primitive operators and their types.
\end{sproof}

\begin{theorem}[Type soundness]
  \label{TypeSoundness}
  If $\typeof{\emptyEnv}{\emptyEnv}{\emptyEnv}{\expr}{\type}$,
  then either
  $\expr$ diverges,
  there exists $\val$ such that $\expr \mapsto^* \val$
  and $\typeof{\emptyEnv}{\emptyEnv}{\emptyEnv}{\val}{\type}$,
  or there exist partial function $\primop$
  and appropriately-typed arguments $\sequence{\val}$
  such that
  $\expr \mapsto^* \ctxt\left[\misapp\right]$.
\end{theorem}
\begin{sproof}[\BODY]
  We argue coinductively using the sequence of reduction steps from $\expr$.
  For any well-typed $\expr$,
  Progress (Lemma \ref{Progress}) implies that either
  $\expr$ has the form $\val$,
  $\expr$ has the form $\ctxt[${\tt ((array () $\primop$) $\val$ $\sequence{}$)}$]$,
  or $\expr \mapsto \expr'$.
  In the first case,
  the reduction sequence terminates in a value,
  so we have $\expr \mapsto^* \val$.
  Furthermore, Preservation (Lemma \ref{Preservation}) implies that
  $\typeof{\emptyEnv}{\emptyEnv}{\emptyEnv}{\val}{\type}$
  In the second case,
  the reduction sequence terminates in a mis-applied primitive operator.
  In the third case,
  Preservation implies that $\typeof{\emptyEnv}{\emptyEnv}{\emptyEnv}{\expr'}{\type}$.
\end{sproof}
